% Generated draft: requires mathematical and source review.
\documentclass[11pt]{article}
\usepackage[margin=25mm]{geometry}
\usepackage{fontspec}
\setmainfont{Latin Modern Roman}
\usepackage{amsmath,amssymb,mathtools,unicode-math}
\setmathfont{Latin Modern Math}
\usepackage{enumitem,needspace,xurl,hyperref}
\setlength{\parindent}{0pt}
\setlength{\parskip}{5pt}
\setlength{\emergencystretch}{3em}
\newcommand{\N}{\mathbb{N}}
\newcommand{\Z}{\mathbb{Z}}
\newcommand{\Q}{\mathbb{Q}}
\newcommand{\R}{\mathbb{R}}
\newcommand{\C}{\mathbb{C}}
\newcommand{\F}{\mathbb{F}}
\newcommand{\A}{\mathbb{A}}
\newcommand{\PP}{\mathbb{P}}
\newcommand{\E}{\mathbb{E}}
\newcommand{\eps}{\varepsilon}
\newcommand{\dd}{\,\mathrm{d}}
\begin{document}
\section*{UnsolvedMath ID: 30000582}
\textbf{Source problem code:} catalogue entry
\subsection{Definitions and mathematical statement}
A Fano manifold is a smooth connected complex projective variety $X$ with ample anticanonical line bundle $K_X^{-1}$. Let $\pi:T_X\to X$ be the total space of its holomorphic tangent bundle, rather than its projectivization. A section is a holomorphic map $s:X\to T_X$ with $\pi\circ s=\operatorname{id}_X$.

Use the source's following rigid-target condition:
\[
\begin{gathered}
\text{For every irreducible projective subvariety }Z\subset T_X
\text{ with }\overline{\pi(Z)}=X,\\
\text{there is a section }s:X\to T_X\text{ such that }Z=s(X).
\end{gathered}
\]
Equivalently in the last line, $\pi|_Z:Z\to X$ is an isomorphism. The classification problem is to determine which Fano manifolds satisfy this condition. Here projective subvariety means a complete projective subvariety inside the total space, not a subvariety of $\PP(T_X)$.

The source also describes rigidity as the property that every deformation of a surjective holomorphic map $f:Y\to X$ arises from an automorphism of $X$. In the fixed-domain, local-family interpretation, a holomorphic deformation $f_t$ with $f_0=f$ should have the form
\[
f_t=a_t\circ f,\qquad a_t\in\operatorname{Aut}(X),\quad a_0=\operatorname{id}_X,
\]
locally in the deformation parameter. The extract does not specify the allowed source spaces $Y$ or whether a local, infinitesimal, or global deformation statement is intended. Its claimed equivalence is retained as source content, while the explicit tangent-bundle condition supplies the main classification target.
\subsection{Short English statement}
Classify Fano manifolds for which every projective subvariety of the tangent-bundle total space that dominates the manifold is a section. The source relates this to rigidity of surjective maps up to target automorphisms, but does not specify the source spaces or deformation conventions.
\subsection{Sources}
\begin{itemize}
\item Supplied report: \emph{Komplexe Analysis} (2007), Oberwolfach Reports, DOI 10.4171/owr/2006/40, pp. 2399--2446; proposed year 2006. \url{https://doi.org/10.4171/owr/2006/40}.
\item Supplied reference labelled ``Fano manifolds and rigid targets'': \url{https://intlpress.com/site/pub/files/_fulltext/journals/pamq/2011/0007/0004/PAMQ-2011-0007-a015.pdf}.
\item Local UnsolvedMath registry, permanent ID 30000582, catalogue entry; created 25 August 2026, updated 16 September 2026. Its corrected\_verified note describes editorial cleanup. The supplied literature triage of 21 August 2026 labels the classification partially solved, reporting a complete Fano-surface classification and broad higher-dimensional sufficient criteria, but no general classification. This dated assessment is retained without independent verification or PDF access.
\end{itemize}
\end{document}
